%% GENERATED FILE -- do not hand-edit.
%% SOURCE: experiments/database/scripts/build_bacteria_candidate_datasets_table.py
\begin{landscape}
\begingroup
\footnotesize
\setlength{\tabcolsep}{3pt}
\renewcommand{\arraystretch}{1.15}
\begin{longtable}{@{}r@{\hspace{3pt}} L{38mm} L{13mm} L{10mm} L{19mm} L{17mm} r@{\hspace{4pt}} r@{\hspace{4pt}} L{19mm} L{74mm}@{}}
\caption[]{The bacterial candidates, ranked by \emph{Meas.} descending, which is
instances times phenotype dimensionality (Sec.~\ref{sec:rule}). \emph{Genotypes} and
\emph{Env} are the perturbation and condition axes. \emph{Inst.} is
genotype$\times$environment records, $\dagger$ where it is the product of the two axes
rather than a reported count and $\ddagger$ where it is an order-of-magnitude estimate.
\emph{Sequence basis} is the route to each strain's total genomic content; a row with no
route is excluded (Table~\ref{tab:bexcluded}). Superscript \textbf{B} marks a row whose
per-record values are not released, \textbf{A} a corpus that re-serves other papers and
is not net new until split by source. Citations and data locations are in
Table~\ref{tab:bsources}; the yeast dataset each row mirrors is in
Table~\ref{tab:banalogs}.}
\label{tab:bfinal}\\
\toprule
\textbf{\#} & \textbf{Dataset (class; phenotype)} & \textbf{Host} & \textbf{Tier} & \textbf{Genotypes} & \textbf{Env} & \textbf{Inst.} & \textbf{Meas.} & \textbf{Sequence basis} & \textbf{Why} \\
\midrule
\endfirsthead
\multicolumn{10}{@{}l}{\footnotesize\emph{Table~\ref{tab:bfinal}, continued}}\\
\toprule
\textbf{\#} & \textbf{Dataset (class; phenotype)} & \textbf{Host} & \textbf{Tier} & \textbf{Genotypes} & \textbf{Env} & \textbf{Inst.} & \textbf{Meas.} & \textbf{Sequence basis} & \textbf{Why} \\
\midrule
\endhead
\bottomrule
\endfoot
1 & \textbf{Fuhrer 2017}\newline {\scriptsize Metabolome / flux; non-targeted flow-injection TOF-MS metabolite ion intensities} & \org{E. coli} & 1 & 3,807 deletions & 1 condition & $3.4\times 10^{4}$ & $2.6\times 10^{8}$ & K-\allowbreak 12-\allowbreak KO & 3,807 Keio single-gene deletion mutants (from 4,320 after excluding 21 very sick strains and 16 injection failures), two independent clones each in technical duplicate, more than 34,000 mass-spectrometric injections. 7,534 distinct m/z features (3,169 negative plus 4,365 positive mode); 3,130 of them were putatively matched to 1,432 of 2,028 chemical formulas, so annotated depth is far below raw feature depth. Deposited in MassIVE and BioStudies, not MetaboLights. \\
\addlinespace[5pt]
2 & \textbf{Wetmore 2015}\newline {\scriptsize Transposon fitness; barcode competitive fitness (BarSeq)} & \org{E. coli} & 1 & 152,018 insertion mutants & 89 conditions & $1.4\times 10^{7}$$\dagger$ & $1.4\times 10^{7}$ & K-\allowbreak 12+\allowbreak transposon & Table 1 gives 152018 strains with unique bar codes for the E. coli library KEIO\_ML9 and fitness estimates for 3471 genes (84 percent). The E. coli lane included 96 samples with 4 time-zero samples across 47 different carbon or nitrogen sources, of which 89 passed quality metrics. Aggregated to gene level the usable record count is 3471 by 89, about 308919. Across all five bacteria the paper reports 387 successful genome-wide assays. No SRA accession is stated in the text; the LBL supplemental URL was taken from the paper and was not fetched. \\
\addlinespace[5pt]
3 & \textbf{Mutalik 2020}\newline {\scriptsize Transposon fitness; barcode competitive fitness (RB-TnSeq BarSeq)} & \org{E. coli} & 1 & 152,018 insertion mutants & 77 conditions & $1.2\times 10^{7}$$\dagger$ & $1.2\times 10^{7}$ & K-\allowbreak 12+\allowbreak transposon & The K-12 BW25113 RB-TnSeq library carries 152018 barcoded insertions in 3728 genes and was assayed against 14 diverse dsDNA phages in 68 RB-TnSeq assays plus 9 no-phage controls, giving the 77 conditions used here. A parallel BL21 library of about 97000 mutants was assayed against 12 phages in 53 pooled fitness experiments plus 9 controls. The same study also ran CRISPRi and Dub-seq arms on the same phage panel, which makes it a three-modality join on one condition axis. Accessions are quoted from the article's data availability statement; the SRA and figshare pages themselves were not fetched. The related Mutalik 2019 Nature Communications Dub-seq paper (doi 10.1038/s41467-018-08177-8, PMID 30659179) is a barcoded overexpression library of 30558 barcode pairs assayed in 155 fitness experiments over 52 chemicals with SRA PRJNA512427 and figshare 10.6084/m9.figshare.6752753.v1; it is omitted as its own row only because plasmid overexpression has no matching modality value here. \\
\addlinespace[5pt]
4 & \textbf{Tong 2020}\newline {\scriptsize Fitness / chemical genomics; 24 hour solid-agar colony growth curve} & \org{E. coli} & 1 & 3,796 deletions & 30 conditions & $1.1\times 10^{5}$ & $8.2\times 10^{6}$ & K-\allowbreak 12-\allowbreak KO & Kinetic growth measurements for 3796 genes under 30 carbon source conditions in MOPS minimal medium, 1536 colonies per plate on solid agar, giving 113880 individual growth curves and over 8 million data points. Plates were scanned every 20 minutes over 24 hours, so each instance is a time series of roughly 72 points; the 72 figure is derived from the stated interval and duration rather than stated directly. 51 poorly annotated genes showed a low growth phenotype in at least one carbon source. This is the cleanest released Keio by carbon-source matrix and the only row here carrying a real time axis. \\
\addlinespace[5pt]
5 & \textbf{PRECISE-1K}\newline {\scriptsize Transcriptome; RNA-seq expression compendium plus 201 iModulon activities} & \org{E. coli} & 3 & 76 designs & 116 conditions & $1.0\times 10^{3}$ & $4.4\times 10^{6}$ & K-\allowbreak 12-\allowbreak KO & 1,035 samples over 4,257 genes, 5 strain backgrounds, 76 unique gene knockouts, 45 distinct projects, yielding 201 iModulons; a combined public K-12 set reaches 2,710 samples and 194 iModulons. The paper reports no single unique-condition total, so 116 is the sum of its itemized condition variables (9 base media, 18 carbon sources, 38 supplements, 42 heterologous proteins, 4 temperatures, 5 pH values). No GEO SuperSeries was found; Zenodo is the stated repository. \\
\addlinespace[5pt]
6 & \textbf{Carruthers 2025}\newline {\scriptsize Production campaign; isoprenol titer by GC-FID plus paired DIA shotgun proteome} & \org{P. putida} & 3 & 472 guides & 1 condition & $1.4\times 10^{3}$$\dagger$ & $2.1\times 10^{6}$ & KT2440+\allowbreak guide & The full text states 472 unique strains (125 single perturbations and 347 combinations) over six DBTL cycles in triplicate, matching the recalled numbers exactly; machine learning selected roughly 400 constructs out of 800,000 possibilities for a fivefold titer gain. ONE study with eight linked deposits: Dryad 10.5061/dryad.gtht76hzh (proteomic metadata plus a 29.7 MB Top3 peptide quantification table, file list enumerated), a Zenodo and GitHub release of per-cycle titer tables (10.5281/zenodo.17178684, JBEI/Isoprenol\_CRISPRi), and seven per-cycle PRIDE accessions PXD063733 (DBTL0), PXD063737, PXD063738, PXD063740, PXD063743, PXD063744, PXD063746 (DBTL6). An ingestion that expects one PXD per paper silently drops six sevenths of the proteomics. \\
\addlinespace[5pt]
7 & \textbf{Lim 2022 putidaPRECISE321}\,\textsuperscript{\textbf{A}}\newline {\scriptsize Transcriptome; batch-normalized log2 TPM RNA-seq expression, decomposed into 84 iModulons} & \org{P. putida} & 4 & 17 wild type & 118 conditions & 321 & $1.8\times 10^{6}$ & reference-\allowbreak only & All recalled numbers confirmed twice over, from the paper and by recomputation on the released matrix: 321 samples, 118 unique condition groups, 21 projects, a 5,564 gene by 321 sample log\_tpm\_norm matrix, 84 iModulons explaining 75.7 percent of variance with 1,265 member genes. The live iModulonDB API returns the dataset under organism p\_putida, dataset precise321. Replicate structure recomputed as 40 conditions at n=2, 71 at n=3, 7 at n=4, summing to 321 exactly. Gene-level values total 1,786,044. This row is an AGGREGATION of 21 projects: 541 samples were collected and 321 passed QC, of which the paper says 305 were previously published and 16 newly generated. The genotype count of 17 distinct genotype tuples (10 with an annotated perturbation such as relA, finR, crc, fleQ) is an UNDERCOUNT, because engineered muconate producers and ALE clones are encoded in condition strings rather than genotype fields. There is no umbrella GEO accession: the sample table carries 47 distinct SRA or BioProject accessions covering 216 of 321 samples and 9 GEO Series, and 105 samples carry no accession at all, so per-sample raw-data provenance is incomplete by construction. \\
\addlinespace[5pt]
8 & \textbf{Borchert 2024 fModules}\,\textsuperscript{\textbf{A}}\newline {\scriptsize Transposon fitness; gene-level RB-TnSeq relative fitness, log2 barcode ratio against the time-zero baseline} & \org{P. putida} & 1 & 4,732 insertion mutants & 179 conditions & $1.6\times 10^{6}$ & $1.6\times 10^{6}$ & KT2440+\allowbreak transposon & Every recalled number checks out, and the 179 versus 183 discrepancy is real rather than an error: Methods say 183 unique growth conditions were collected, the Abstract says 179 were analyzed, and the released workbook's classifier column has exactly 179 distinct values, so 183 is the collection and 179 is the analyzed matrix. Fitness for 4,732 of 5,564 protein-coding genes. The instance count is VERIFIED, not a product: fModule\_Metadata.xlsx (23.3 MB) was opened and its fitness sheet counted at 4,732 rows by 332 value columns with zero nulls, alongside a t-statistic sheet of identical shape. This row is an AGGREGATION and a superset of the other transposon rows: 254 of 332 samples were pulled from the Fitness Browser and 78 are the Beckham-group sets (Putida\_ML5\_set100 and set101) deposited as SRA PRJNA809672 (15 runs), PRJNA856070 (57 runs) and PRJNA1011287 (39 runs), all counted. Replicate structure: 194 samples from duplicate conditions, 78 from triplicate, 60 singletons. The environment axis is two-slot (condition\_1 plus condition\_2 with units), which is what keeps glucose-plus-stressor tolerance samples distinguishable from sole-carbon catabolism samples; collapsing it to one string destroys that contrast. Caveat: genes lacking a value in any sample were dropped, so condition-specific genes reported by the primary papers are silently absent here. \\
\addlinespace[5pt]
9 & \textbf{Caglar 2017}\newline {\scriptsize Multi-omics campaign; matched RNA-seq and shotgun proteomics plus central carbon fluxes} & \org{E. coli} & 4 & 1 wild type & 34 conditions & 322 & $1.4\times 10^{6}$ & reference-\allowbreak only & 34 growth conditions in E. coli B REL606 with 152 RNA-seq samples, 105 proteomics samples and 65 flux measurements (322 total), spanning exponential and stationary phase plus two time courses from 3 hours to 2 weeks. The paper reports 4,196 distinct mRNAs and proteins as a joint figure, so treat 4,196 as the matched transcript-and-protein gene axis rather than as two independent depths. \\
\addlinespace[5pt]
10 & \textbf{Nichols 2011}\newline {\scriptsize Fitness / chemical genomics; colony size drug-gene fitness score} & \org{E. coli} & 1 & 3,979 deletions & 324 conditions & $1.3\times 10^{6}$$\dagger$ & $1.3\times 10^{6}$ & K-\allowbreak 12-\allowbreak KO & The paper states 3979 mutant strains passing quality control, 324 conditions covering 114 unique stresses, and a total of 13497 significant phenotypes at 5 percent FDR, which is about 1 percent of all condition-gene pairs tested. The URL printed in the paper, http://ecoliwiki.net/tools/chemgen/, now returns HTTP 410 Gone; the same portal is live at the ecoliwiki.org host and serves the full per-strain by per-condition flat file. The Keio background is BW25113, not MG1655 itself. \\
\addlinespace[5pt]
11 & \textbf{Goodall 2018}\newline {\scriptsize Transposon fitness; TraDIS insertion-site read density, essentiality call} & \org{E. coli} & 1 & 901,383 insertion mutants & 1 condition & $9.0\times 10^{5}$$\dagger$ & $9.0\times 10^{5}$ & K-\allowbreak 12+\allowbreak transposon & A transposon mutant library estimated at approximately 3.7 million mutants, sequenced to 8279309 mapped reads over 901383 unique insertion sites, from which 358 genes were identified as essential. Only one condition was used, Luria broth at 37 degrees C to OD600 of 1.0, in two independent DNA extracts TL1 and TL2. 248 genes (59.5 percent) of the essential calls were common to this study, the Keio collection and PEC. The ENA accession is quoted from the paper's data availability statement and the ENA page itself was not fetched. \\
\addlinespace[5pt]
12 & \textbf{Price 2018}\newline {\scriptsize Transposon fitness; barcode competitive fitness (BarSeq log ratio)} & \org{E. coli} & 1 & 3,789 insertion mutants & 162 conditions & $6.1\times 10^{5}$$\dagger$ & $6.1\times 10^{5}$ & K-\allowbreak 12+\allowbreak transposon & The E. coli subset is separable: the published per-organism page for orgId Keio (Escherichia coli BW25113) states 207 condition samples with 162 successful, and its fit\_logratios\_good.tab was downloaded and counted at 162 experiment columns by 3789 gene rows. Whole-study totals are 32 bacteria, 4870 genome-wide fitness experiments meeting consistency criteria, and 11779 poorly annotated protein-coding genes with phenotypes. A 2021 data correction removes the sucrose and D-mannitol experiments for this organism. The live Fitness Browser sits behind a Cloudflare challenge and could not be read, so its current counts may exceed these. \\
\addlinespace[5pt]
13 & \textbf{Choe 2025}\newline {\scriptsize CRISPR library screen; sgRNA abundance by next-generation sequencing} & \org{E. coli} & 2 & 39,591 guides & 12 conditions & $4.8\times 10^{5}$$\dagger$ & $4.8\times 10^{5}$ & K-\allowbreak 12+\allowbreak guide & The initial library holds 39591 sgRNA sequences at 99.96 percent coverage, targeting 4198 coding sequences. Twelve chemical conditions were evaluated: CCCP, polymyxin B, pyocyanin, rifampicin, sulfamethizole, verapamil, puromycin, erythromycin, phleomycin, mitomycin C, MMS and novobiocin. This is the genuine antibiotic-panel CRISPRi screen in E. coli and is the correct source for that claim rather than Rousset 2018. The ENA project was fetched and confirmed, released December 2024, submitted by KAIST. \\
\addlinespace[5pt]
14 & \textbf{Rapp 2026 CRISPRi metabolome}\newline {\scriptsize CRISPR library screen; flow-injection MS metabolome (ion intensities), plus carotenoid titer in the engineering follow-up} & \org{E. coli} & 3 & 1,515 guides & 1 condition & $1.5\times 10^{3}$ & $4.5\times 10^{5}$ & K-\allowbreak 12+\allowbreak guide & Arrayed CRISPRi library covering all 1515 metabolic genes of iML1515, each strain profiled by fast flow-injection MS, so per-design metabolome vectors are recoverable; 36\% of iML1515-predicted metabolites accumulated specifically in some knockdown, and LC-MS/MS spectra were generated for 102 previously uncharacterized metabolites. The exact number of annotated ions per strain was not verified from the full text, so dimensionality 300 is an estimate. The MassIVE record was fetched and confirmed. \\
\addlinespace[5pt]
15 & \textbf{Schmidt 2022 nitrogen}\newline {\scriptsize Transposon fitness; gene-level RB-TnSeq fitness across nitrogen-containing sole nitrogen sources} & \org{P. putida} & 1 & 4,778 insertion mutants & 71 conditions & $3.4\times 10^{5}$$\dagger$ & $3.4\times 10^{5}$ & KT2440+\allowbreak transposon & Not on the requested list but it clearly belongs, and it is the largest single-study condition axis for this organism: assimilation of 52 different nitrogen-containing compounds across 71 tested conditions, with significant fitness phenotypes in 672 genes including 100 transcriptional regulators and 112 transport proteins. Same JBEI-1 library and same Deutschbauer and Keasling pipeline as the lysine, valerolactam and fatty-acid rows, and it is one of the source studies cited by the Borchert 2024 compendium, so it must be deduplicated against that row. The compendium's nitrogen-source experiment group holds 104 of the 332 samples, which is where most of these land. Fitness Browser only, Cloudflare-blocked, so unconfirmed. \\
\addlinespace[5pt]
16 & \textbf{Wang 2018}\newline {\scriptsize CRISPR library screen; sgRNA abundance by Illumina sequencing, gene and sgRNA fitness} & \org{E. coli} & 2 & 56,071 guides & 5 conditions & $2.8\times 10^{5}$$\dagger$ & $2.8\times 10^{5}$ & K-\allowbreak 12+\allowbreak guide & Two corrections to the commonly cited form of this row. The venue is Nature Communications 9:2475, not Science; no Science 2018 genome-wide E. coli CRISPRi paper could be found. The library is 55671 targeting sgRNAs plus 400 negative controls, not about 92000; the 92000-guide figure belongs to the Bikard-lab libraries of Cui 2018 and Rousset 2018. Coverage is at least one sgRNA for 98.6 percent of 4140 protein-coding genes and 79.8 percent of 178 RNA-coding genes. The five screened phenotypes are essentiality in LB, auxotrophy in MOPS versus LB, L-tryptophan biosynthesis, furfural tolerance and isobutanol tolerance. Processed per-gene and per-sgRNA fitness are released as supplementary data; no raw-sequencing accession could be confirmed. \\
\addlinespace[5pt]
17 & \textbf{Rousset 2018}\newline {\scriptsize CRISPR library screen; guide abundance log2 fold change (relative sgRNA fitness)} & \org{E. coli} & 2 & 59,000 guides & 4 conditions & $2.4\times 10^{5}$$\dagger$ & $2.4\times 10^{5}$ & K-\allowbreak 12+\allowbreak guide & The screen starts from a pool of about 92000 sgRNAs targeting random chromosomal positions and is filtered to about 59000 guides. Conditions are growth in rich medium over 17 generations for essentiality plus infection by phages lambda, T4 and 186 at MOI 1, all in triplicate; the paper screens no antibiotics, so any title citing antibiotic exposure for this paper is a conflation. 379 genes were selected for follow-up, including 235 annotated as essential. The ENA accession was verified directly at EBI. \\
\addlinespace[5pt]
18 & \textbf{Shiver 2016}\newline {\scriptsize Fitness / chemical genomics; colony opacity fitness score (Iris)} & \org{E. coli} & 1 & 3,975 deletions & 57 conditions & $2.3\times 10^{5}$$\dagger$ & $2.3\times 10^{5}$ & K-\allowbreak 12-\allowbreak KO & 3975 Keio deletion mutants against 26 new stresses for 57 conditions in total, scored from colony opacity with the Iris image analysis software, and designed as an extension of the Nichols 2011 condition set. The release is raw plate images plus Iris output and ranked score files rather than a clean gene by condition matrix, so building the matrix requires rerunning the published scoring scripts. The Zenodo mirror of the Dryad deposit was fetched and its file list confirmed. \\
\addlinespace[5pt]
19 & \textbf{Thompson 2020 fatty acid and alcohol}\newline {\scriptsize Transposon fitness; gene-level RB-TnSeq fitness on fatty acid and alcohol carbon sources} & \org{P. putida} & 1 & 4,778 insertion mutants & 23 conditions & $2.2\times 10^{5}$$\dagger$ & $2.2\times 10^{5}$ & KT2440+\allowbreak transposon & The recalled 13 fatty acids and 10 alcohols is confirmed verbatim in the abstract and the compounds were enumerated: straight-chain C3 to C10 plus C12 and C14, the esters Tween 20 and butyl stearate, oleic acid; and ethanol, butanol, pentanol, 1,2-propanediol, 1,3-butanediol, 1,4-butanediol, 1,5-pentanediol, isopentanol, isoprenol, 2-methyl-1-butanol. Biological duplicate throughout, corroborated by the compendium metadata showing exactly 2 samples per compound, hence 23 x 2 x 4,778; the condition-collapsed count is about 109,900. Directly relevant to the isoprenol rows: this paper postulates the catabolic routes for isoprenol and isopentanol through leucine metabolism after oxidation and CoA activation, which is the product-loss pathway the isoprenol producers delete. A published correction exists (AEM 87(8):e00177-21, DOI 10.1128/aem.00177-21). Fitness Browser only, Cloudflare-blocked, so unconfirmed; not every published condition survived the compendium's completeness filter, so the full matrix is not recoverable from the workbook alone. \\
\addlinespace[5pt]
20 & \textbf{Borchert 2023 lignin tolerance}\newline {\scriptsize Transposon fitness; gene-level RB-TnSeq fitness under lignin-stream inhibitors on a fixed glucose background} & \org{P. putida} & 1 & 4,732 insertion mutants & 16 conditions & $2.1\times 10^{5}$$\dagger$ & $2.1\times 10^{5}$ & KT2440+\allowbreak transposon & This is the lignin-stream tolerance row with reconstructed engineering strategies, and it is a tolerance rather than a catabolism screen: every stressor is overlaid on M9 plus 20 mM glucose. The 15 stressors were enumerated from the SRA sample titles (ferulate, 4-coumarate, levulinate, beta-ketoadipate, protocatechuate, cis,cis-muconate, glycolate, lactate, acetate, Na2SO4, NaCl, vanillate, vanillin, 4-hydroxybenzoate, 4-hydroxybenzaldehyde) plus a glucose-only control, in triplicate: the 57 SRA runs decompose exactly as 15 x 3 selective, plus 3 plus 3 glucose-control enrichments, plus 3 plus 3 time-zero samples. PRJNA856070 confirmed by fetching the BioProject page (57 SRA experiments, 57 BioSamples). Reconstructed strategies from the fitness hits: delta-gacAS, delta-fleQ, delta-lapAB, delta-ttgR with Ptac-ttgABC, Ptac-PP\_1150-PP\_1152, delta-relA, delta-PP\_1430, several of which also improved growth in a complex lignin-stream mimic. These samples are NOT in the Fitness Browser, which is precisely why the Borchert 2024 compendium reached into SRA for them; the 39 samples of Putida\_ML5\_set100 are recoverable at full precision from that workbook. \\
\addlinespace[5pt]
\midrule \multicolumn{10}{@{}l}{\textbf{End of tranche 1. Rows 1--20 are the first builds, 6 per host.}}\\ \midrule
21 & \textbf{Schastnaya 2021}\newline {\scriptsize Metabolome / flux; metabolome of phosphosite point mutants and matched deletion mutants} & \org{E. coli} & 4 & 89 designs & 5 conditions & 445$\dagger$ & $2.0\times 10^{5}$ & engineered-\allowbreak chassis & 52 known phosphosites on 23 central metabolic enzymes were mutated in E. coli MG1655 delta-mutS, giving 89 phosphomutants profiled on 5 carbon sources against matched Keio deletion strains, with 460 to 500 annotated metabolites by flow-injection TOF-MS. This is the best genotype-by-condition phospho-regulation panel, but it generated no phosphoproteome of its own; its phosphosite list is drawn from earlier studies such as Potel 2018. \\
\addlinespace[5pt]
22 & \textbf{Cui 2018}\newline {\scriptsize CRISPR library screen; guide abundance log2 fold change over 17 generations} & \org{E. coli} & 2 & 92,000 guides & 2 conditions & $1.8\times 10^{5}$$\dagger$ & $1.8\times 10^{5}$ & K-\allowbreak 12+\allowbreak guide & The library is about 92000 unique guide RNAs targeting random positions along the MG1655 genome with an NGG PAM requirement, averaging 19 targets per gene; an exact targeted-gene count is not stated. The two genome-wide screens differ only in dCas9 expression level (strains LC-E18 and LC-E75), both in rich medium with anhydrotetracycline in triplicate, with no stress panel. Data availability names Supplementary Data 4 for the screen results and otherwise defers to the corresponding author on request, so there is no SRA, ENA or GEO accession for this paper. \\
\addlinespace[5pt]
23 & \textbf{Schmidt 2016}\newline {\scriptsize Proteome; absolute protein copies per cell and per cell volume} & \org{E. coli} & 4 & 3 wild type & 22 conditions & 66$\dagger$ & $1.6\times 10^{5}$ & reference-\allowbreak only & 2,359 proteins quantified across all 22 conditions, about 55 percent of predicted ORFs, calibrated to copies per cell with flow cytometry and synthetic peptide standards. Primary strain BW25113, with MG1655 and NCM3722 at selected conditions; biological triplicates for the main dataset, so 66 is a product not a reported sample count. \\
\addlinespace[5pt]
24 & \textbf{Butland 2008}\,\textsuperscript{\textbf{B}}\newline {\scriptsize Genetic interaction; colony size interaction (S) score} & \org{E. coli} & 1 & 155,415 double mutants & 1 condition & $1.6\times 10^{5}$$\ddagger$ & $1.6\times 10^{5}$ & K-\allowbreak 12-\allowbreak KO x KO & Only the citation fields and the abstract-level method description could be confirmed: a quantitative screening procedure based on conjugation of E. coli deletion or hypomorphic strains to create double mutants on a genome-wide scale. The article is closed access with no PMC deposit, and nature.com, Springer static content and mirror sites all refused fetching, so the number of query genes, array strains and screened pairs are all unconfirmed. A search snippet suggesting 39 genome-wide screens, which would give roughly 39 by 3985 or 155415 pairs, could not be verified on any fetched page, so the counts above are an unverified estimate and the row is marked blocked pending a paywalled retrieval. \\
\addlinespace[5pt]
25 & \textbf{Mori 2021}\newline {\scriptsize Proteome; absolute protein mass fractions by DIA/SWATH with xTop inference} & \org{E. coli} & 4 & 3 wild type & 60 conditions & 66 & $1.5\times 10^{5}$ & reference-\allowbreak only & 2,335 proteins with absolute mass fractions over about 60 growth conditions in 66 samples, covering carbon, nitrogen, phosphate and oxygen limitation, non-metabolic stresses and biofilm states, in K-12 NCM3722 and derivatives plus MG1655 sub-strain EQ353 and E. coli Nissle 1917. No ProteomeXchange or PRIDE accession was stated in the text examined; the deposit found is the SWATHAtlas spectral library. An author correction was published in October 2024 and should be read with the paper. \\
\addlinespace[5pt]
26 & \textbf{Ishii 2007}\newline {\scriptsize Multi-omics campaign; paired transcriptome, proteome, metabolome and 13C flux in glucose-limited chemostats} & \org{E. coli} & 4 & 25 designs & 5 conditions & 29 & $1.2\times 10^{5}$ & K-\allowbreak 12-\allowbreak KO & Design confirmed from a secondary open-access source: 24 single-gene disruptants at a fixed dilution rate of 0.2 per hour plus wild type at 5 dilution rates, measured by DNA microarray, 2D-DIGE, CE-TOFMS and metabolic flux analysis. Per-layer depth (transcripts, proteins, metabolites, fluxes) is NOT confirmed: the paper is paywalled and absent from PMC, and no accessible secondary source quotes those counts, so 4,300 is a genome-wide-array estimate for the dominant layer and must be replaced before use in any ranking. \\
\addlinespace[5pt]
27 & \textbf{Campos 2018}\newline {\scriptsize Fitness / chemical genomics; single-cell microscopy morphology and cell cycle features} & \org{E. coli} & 2 & 4,227 deletions & 1 condition & $4.2\times 10^{3}$$\dagger$ & $1.1\times 10^{5}$ & K-\allowbreak 12-\allowbreak KO & 4227 strains of the Keio collection, covering 98 percent of the non-essential genome and 87 percent of the complete genome, imaged in one condition (M9 with 0.1 percent casamino acids and 0.2 percent glucose at 30 degrees C). 26 quantitative features per strain: 19 morphological, 5 cell cycle, and 2 growth-related. On average about 360 cells were imaged per strain, roughly 1.3 million cells after filtering. This is the highest-dimensionality E. coli deletion phenotype release found, and the natural counterpart to a yeast morphology or transcriptome profile rather than to a scalar fitness score. \\
\addlinespace[5pt]
28 & \textbf{Typas 2008}\,\textsuperscript{\textbf{B}}\newline {\scriptsize Genetic interaction; colony pixel count interaction score} & \org{E. coli} & 1 & 91,655 double mutants & 2 conditions & $9.2\times 10^{4}$$\dagger$ & $9.2\times 10^{4}$ & K-\allowbreak 12-\allowbreak KO x KO & This is a method paper (GIANT-coli). Its released per-pair data cover only two demonstration screens: a 12 by 12 cross yielding 66 distinct pairwise double mutants plus 12 self-matings, and the pal and yraP screens in Supplementary Tables 2A and 2B. The genome-wide work is reported only as a linkage curve from 9 by 3985 crosses in LB and 14 by 3985 in M9, which is the 91655 figure above and is a product, not a released matrix. Status is blocked because no genome-wide double-mutant score matrix was released. The toolkit couples the Keio kanamycin deletion library with the ASKA chloramphenicol library, about 4000 single-gene deletions each. \\
\addlinespace[5pt]
29 & \textbf{Girgis 2009}\newline {\scriptsize Transposon fitness; microarray genetic footprinting z-score per locus} & \org{E. coli} & 1 & 500,000 insertion mutants & 17 conditions & $7.3\times 10^{4}$$\ddagger$ & $7.3\times 10^{4}$ & K-\allowbreak 12+\allowbreak transposon & The library is about 5 by 10\textasciicircum{}5 mutants each with a single transposon insertion, selected in 17 antibiotics, read out by hybridizing transposon-junction DNA against genomic DNA on spotted microarrays. The released data are per-locus z-scores, with Dataset S5 the combined z-scores across all loci, so usable records are locus by drug rather than strain by drug; the number of array loci is not stated, so the instance count here is an estimate from an assumed roughly 4300 loci and is labeled as such. No GEO or ArrayExpress accession exists for this dataset; a GEO search returned only the unrelated expression series GSE10855. The companion motility screen is Girgis HS, Liu Y, Ryu WS, Tavazoie S 2007 PLoS Genetics 3(9):e154, doi 10.1371/journal.pgen.0030154, PMID 17941710, same library size, four selection regimes released as 13 supplementary microarray datasets with no accession. \\
\addlinespace[5pt]
30 & \textbf{Banerjee 2025}\newline {\scriptsize Production campaign; indigoidine titer from p-coumarate plus growth rate, with a global DIA proteome} & \org{P. putida} & 4 & 12 promoter variants & 3 conditions & 36$\ddagger$ & $7.2\times 10^{4}$ & KT2440+\allowbreak promoter & This is the recalled genome-scale design-tradeoff campaign, and it is the cleanest promoter-variant genotype axis I found for KT2440: a four-gene growth-coupling cutset for p-coumarate to glutamine to indigoidine, whose fourth gene PP\_0897 (a fumarate hydratase isomer) proved multifunctional and rate-limiting, so its native promoter was swapped for the Anderson-collection pJ23109 or the weak PP\_0415 promoter. PXD050285 verified two ways (ProteomeCentral and the PRIDE v3 API) and its files enumerated: 26 raw MS runs, Orbitrap Exploris 480, DIA searched library-free with DIA-NN at 1 percent global FDR, comprising an 8-run promoter-variant arm (pJ and 0415 driving PP\_0897, 4 replicates each) and an 18-run cross-feeding arm (strains 2370 and 2487 across M9 alanine, M9 alanine malate, M9 p-CA alanine malate, n=3). The about 2,000 quantified proteins is REPORTED verbatim in the paper, not estimated. A separate 190-substrate BIOLOG grid in the same paper is a much larger environment axis if it is ingested. Note the PRIDE submission title differs from the published title; both are correct for their object. \\
\addlinespace[5pt]
31 & \textbf{Fang 2025 FFA CRISPRi-FACS}\,\textsuperscript{\textbf{B}}\newline {\scriptsize CRISPR library screen; sgRNA enrichment (FACS sort plus NGS), then titer (g/L) on validated strains} & \org{E. coli} & 2 & 55,671 guides & 1 condition & $5.6\times 10^{4}$ & $5.6\times 10^{4}$ & K-\allowbreak 12+\allowbreak guide & The screen reuses a previously published plasmid library of 55671 sgRNAs, sorted with a fluorescent free-fatty-acid biosensor and read out by NGS; the engineered pcnBi-acrDi-fadR+ strain reached 35.1 g/L FFAs in fed-batch. Pooled enrichment gives a guide-level score rather than a per-design titer, so only the handful of individually reconstructed strains carry phenotype values; no sequencing accession was confirmed. \\
\addlinespace[5pt]
32 & \textbf{Babu 2014}\newline {\scriptsize Genetic interaction; colony size epistasis (S) score} & \org{E. coli} & 1 & 671,071 double mutants & 1 condition & $4.3\times 10^{4}$ & $4.3\times 10^{4}$ & K-\allowbreak 12-\allowbreak KO x KO & The largest E. coli digenic screen that could be verified. It reports over 600000 digenic mutant combinations from 163 query donor genes crossed into an array of 3968 non-essential deletions plus 149 hypomorphic strains, which is 671071 nominal pairs, with donors transferred by conjugation from an Hfr-Cavalli chloramphenicol-marked background. Only the high-confidence subset is released, 25239 aggravating plus 17466 alleviating for 42705 pairs in Table S2, so the full matrix is not recoverable and the instance count reflects the released subset. The dedicated portal http://ecoli.med.utoronto.ca/esga no longer resolves to the Emili lab (the host presents a certificate for lymnaea.org), and BioGRID has no entry for PMID 24586182. The earlier Babu 2011 PLoS Genetics map (doi 10.1371/journal.pgen.1002377, PMID 22125496) is smaller at over 235000 combinations and likewise released only significant pairs. Kumar 2016 Cell Reports (doi 10.1016/j.celrep.2015.12.060, PMID 26774489) covers 102255 viable digenic mutants and is the one E. coli interaction set BioGRID serves, at 31644 curated interactions. \\
\addlinespace[5pt]
33 & \textbf{Gupta 2024 turnover}\newline {\scriptsize Translation / turnover; per-protein degradation and turnover rates} & \org{E. coli} & 3 & 1 wild type & 13 conditions & 13$\ddagger$ & $4.2\times 10^{4}$ & reference-\allowbreak only & About 3,200 proteins, 77 percent of all genes, with turnover rates across 13 growth conditions in strain NCM3722, by heavy 15N dynamic labeling with TMTproC complement-reporter mass spectrometry. Replicates exist but the per-condition replicate count is unconfirmed, so instances is a floor. Cite the 2024 Nature Communications version, not the 2022 bioRxiv preprint. \\
\addlinespace[5pt]
34 & \textbf{MCF2Chem 2023}\,\textsuperscript{\textbf{A}}\newline {\scriptsize Aggregation / support; titer, yield, productivity, content per literature record} & \org{E. coli} & 3 & 8,888 designs & 1 condition & $8.9\times 10^{3}$ & $3.6\times 10^{4}$ & reference-\allowbreak only & 8888 production records covering 1231 compounds in 590 microbial cell factories, with titer, yield, productivity and content plus strain culture information; bacteria are 60\% of species, and E. coli, S. cerevisiae, Y. lipolytica and C. glutamicum together account for 78\% of the compounds, with E. coli alone producing about a quarter of the compounds among the top 20 species. No exact E. coli record count is published, so the E. coli subset size could not be determined. These are counts of literature records, not measured strain-condition observations, and must be resolved to primary source DOIs before they count as independent datasets. \\
\addlinespace[5pt]
35 & \textbf{CeCaFDB flux compendium}\,\textsuperscript{\textbf{A}}\newline {\scriptsize Aggregation / support; curated published 13C flux distributions on a normalized central-carbon network} & \org{E. coli} & 3 & mixed & not stated & 297 & $2.3\times 10^{4}$ & reference-\allowbreak only & 581 flux distributions across 36 organisms from 118 references (1995 to 2013); E. coli alone contributes 297 flux distributions from 32 references, the largest single-organism set (S. cerevisiae is next at 76). The normalized comparison network has 66 metabolites and 76 reactions, which is the per-map dimensionality bound, not necessarily each source study's own network size. The site returned an expired-TLS error on 2026-09-24, so current operation is unconfirmed. \\
\addlinespace[5pt]
36 & \textbf{Thompson 2019 lysine}\newline {\scriptsize Transposon fitness; gene-level RB-TnSeq fitness on glucose, 5-aminovalerate, D-lysine and L-lysine} & \org{P. putida} & 1 & 4,778 insertion mutants & 4 conditions & $1.9\times 10^{4}$$\dagger$ & $1.9\times 10^{4}$ & KT2440+\allowbreak transposon & This is where the name JBEI-1 is defined, verbatim: the JBEI-1 library was created by diluting a 1 mL aliquot of the previously described P. putida RB-TnSeq library (Rand 2017) into 500 mL of LB plus kanamycin. So JBEI-1 is a regrown working stock of Putida\_ML5, not a separate library, which is why every transposon row here shares one join key. Four sole-carbon conditions at 10 mM in MOPS minimal medium, 48 h, with three time-zero aliquots; Methods state no per-condition replicate count. The paper reports 39 genes with fitness below -2 at absolute t above 4. Fitness data are said to be publicly available at fit.genomics.lbl.gov with no per-experiment accession, and that site returns a Cloudflare challenge, so unconfirmed. Strains and plasmids at public-registry.jbei.org folder 391. \\
\addlinespace[5pt]
37 & \textbf{Rand 2017 Putida\_ML5 library}\newline {\scriptsize Modality / backbone; defines the barcoded insertion pool every fitness row is scored on; itself assayed on four carbon sources} & \org{P. putida} & 3 & 185,401 insertion mutants & 4 conditions & $1.9\times 10^{4}$$\ddagger$ & $1.9\times 10^{4}$ & KT2440+\allowbreak transposon & CORRECTION to the recalled library definition on two points. First, the library is named Putida\_ML5 and was built here in the Deutschbauer and Arkin lab from the pKMW3 mariner vector library of Wetmore 2015; JBEI-1 is not a separate library but a regrown, re-aliquoted working stock of Putida\_ML5 named later in Thompson 2019 mBio. The compendium metadata's mutantLibrary column carries exactly two values, Putida\_ML5 (25 samples) and Putida\_ML5\_JBEI (307 samples). Second, the recalled sizes of about 100,000 insertion mutants and about 4,800 nonessential genes are BOTH wrong. The only quantitative description in the literature is in Borchert 2022 (ACS Synth Biol 11:2015-2021, DOI 10.1021/acssynbio.2c00119, PMID 35657709) citing this paper: 185,401 uniquely barcoded transposon insertions, of which 32,591 are intergenic and 152,810 map to 5,213 of 5,661 genes. The about 4,800 figure is a FITNESS-coverage number, not library coverage: the Fitness Browser states fitness data for 4,778 of 5,661 genes, and Borchert 2024 reports 4,732 of 5,564 after filtering. Rand 2017 itself says only thousands of kanamycin-resistant colonies. Its own screen covers 4 carbon sources in duplicate over two days. Join on barcode to mapped insertion to PP\_ locus tag against AE015451 (6,181,873 nt, taxid 160488); no study releases per-barcode fitness, all publish gene-level aggregates. The Fitness Browser organism page is orgId=Putida and reported 314 experiments across 54 carbon sources, 52 nitrogen sources, 2 stress compounds and 27 other conditions, but it is behind a Cloudflare managed JavaScript challenge today, so those figures were read from Internet Archive snapshots (2025-03-10 and 2025-04-24) and must NOT be recorded as live 2026 values. Hence accession\_confirmed false. \\
\addlinespace[5pt]
38 & \textbf{Balakrishnan 2022}\newline {\scriptsize Multi-omics campaign; matched mRNA and protein mass fractions plus promoter activity and mRNA degradation rates} & \org{E. coli} & 4 & 2 designs & 4 conditions & 8$\dagger$ & $1.5\times 10^{4}$ & reference-\allowbreak only & The cleanest transcript-and-protein join on identical samples: more than 1,900 proteins with matched mRNA, plus about 2,700 genome-wide mRNA degradation rates, in E. coli K-12 NCM3722 (and a delta-rsd mutant) across a reference glucose minimal medium plus carbon limitation, anabolic limitation and translational inhibition, growth rates 0.3 to 0.9 per hour. Replicate counts vary by assay and are not uniformly reported, so instances is a product. \\
\addlinespace[5pt]
39 & \textbf{Brunk 2016}\newline {\scriptsize Multi-omics campaign; 86 metabolites plus 55 protein complexes, with product titers (isopentenol, limonene, bisabolene)} & \org{E. coli} & 4 & 9 pathway designs & 9 conditions & 81$\dagger$ & $1.1\times 10^{4}$ & engineered-\allowbreak chassis & Nine strains (eight engineered plus wild-type E. coli DH1) carrying three versions of a heterologous mevalonate pathway, sampled across a 0 to 72 hour batch fermentation: the joint panel is 86 metabolites and 55 protein complexes at 9 time points, while the aggregate metabolomics set is described as 9 strains x 13 time points x 86 metabolites, so instances 81 is the product for the joint panel and the metabolome alone is larger. Variances come from triplicate measurements where available. No omics repository accession was confirmed. \\
\addlinespace[5pt]
40 & \textbf{Thompson 2019 valerolactam}\newline {\scriptsize Transposon fitness; gene-level RB-TnSeq fitness on lactam carbon sources, then valerolactam titer in deletion strains} & \org{P. putida} & 2 & 4,778 insertion mutants & 2 conditions & $9.6\times 10^{3}$$\dagger$ & $9.6\times 10^{3}$ & KT2440+\allowbreak transposon & Two selective RB-TnSeq conditions, 10 mM valerolactam and 10 mM 5-aminovalerate in MOPS minimal medium, with three time-zero aliquots; the two 500 uL aliquots per condition are pooling, not replication, so the paper reports no replicate count. RB-TnSeq is the small half of this paper; the payload worth ingesting is the engineering ladder it produced, wild type 0 mg/L to delta-oplBA 9.27 to delta-oplBA delta-davT 85.19 to delta-oplBA delta-davT delta-alr 91.97 mg/L valerolactam at 48 h, with strains at public-registry.jbei.org folder 456. Fitness lives only in the Fitness Browser with no per-experiment accession, and that site is behind a Cloudflare managed challenge to both WebFetch and curl, so accession\_confirmed is false. The two valerolactam carbon-source samples are however confirmed to exist, appearing by name (set6IT064, set7IT045) in the Borchert 2024 compendium workbook. \\
\addlinespace[5pt]
41 & \textbf{D2Cell 2026}\,\textsuperscript{\textbf{A}}\newline {\scriptsize Aggregation / support; literature-reported production metric per gene-modification record} & \org{E. coli} & 3 & 8,134 designs & 1 condition & $8.1\times 10^{3}$ & $8.1\times 10^{3}$ & reference-\allowbreak only & Confirmed verbatim in the preprint: 8134 experimentally reported single- and double-gene modification records for 73 products in E. coli, merged for training with 11643 single-gene-modification rows for the same 73 products simulated by the genome-scale model (FSEOF-style), so the simulated rows originate from a separate generation step and are separable in principle, but the preprint states no explicit flag distinguishing them in the released file; the whole database is 29006 entries, 1210 products, 751 organisms. These are counts of literature records, not measured strain-condition observations, and must be resolved to primary source DOIs before they count as independent datasets. \\
\addlinespace[5pt]
42 & \textbf{Li 2014}\newline {\scriptsize Translation / turnover; absolute protein synthesis rates from ribosome profiling} & \org{E. coli} & 3 & 1 wild type & 2 conditions & 2$\dagger$ & $6.1\times 10^{3}$ & reference-\allowbreak only & 3,041 genes accounting for more than 96 percent of total protein synthesized in rich defined medium, with a similar number in glucose minimal medium, at 90 million fragments per sample. Method is ribosome profiling calibrated to total protein per doubling; the quantitative mass spectrometry comparison is to published external datasets, not generated here. Strain designation is unconfirmed (main-text methods do not name it). \\
\addlinespace[5pt]
43 & \textbf{Potel 2018 phosphoproteome}\newline {\scriptsize Proteome; phosphosite identifications including phosphohistidine} & \org{E. coli} & 4 & 1 wild type & 2 conditions & 2$\ddagger$ & $4.3\times 10^{3}$ & reference-\allowbreak only & The deepest E. coli phosphoproteome found, but with almost no breadth: one wild-type MG1655 strain on glycerol, exponential versus stationary phase. The 2,129 total phosphosites and 246 phosphohistidine sites on 173 proteins come from a secondary search snippet, not from a fetched primary text; the abstract confirms only that about 10 percent of sites are phosphohistidine, so dimensionality needs primary confirmation before ranking. \\
\addlinespace[5pt]
44 & \textbf{Baba 2006}\newline {\scriptsize Modality / backbone; deletion viability (in-frame knockout obtained or not)} & \org{E. coli} & 3 & 3,985 deletions & 1 condition & $4.0\times 10^{3}$ & $4.0\times 10^{3}$ & K-\allowbreak 12-\allowbreak KO & Of 4288 genes targeted, mutants were obtained for 3985, held in duplicate for 7970 strains; 303 genes including 37 of unknown function could not be disrupted and are the paper's essential-gene candidates. Strain background is E. coli K-12 BW25113. The distribution URL given in the paper (GenoBase, http://ecoli.aist-nara.ac.jp/) is stale; NBRP E. coli currently lists 3909 Keio entries ready to distribute. This is the strain resource and the join axis for Nichols 2011, Tong 2020, Campos 2018, Shiver 2016 and the conjugation-based interaction screens, not a phenotype screen itself. \\
\addlinespace[5pt]
45 & \textbf{Gerdes 2003}\newline {\scriptsize Transposon fitness; genetic footprinting essentiality assertion} & \org{E. coli} & 2 & 3,746 insertion mutants & 1 condition & $3.7\times 10^{3}$ & $3.7\times 10^{3}$ & K-\allowbreak 12+\allowbreak transposon & Unambiguous essentiality assessments were made for 3746 (87 percent) of E. coli protein-encoding genes or ORFs, of which 620 (14 percent) were asserted essential and 3126 (73 percent) non-essential; no assertion was possible for 327 genes on technical grounds and evidence was insufficient for 218. Method is Tn5-based genetic footprinting in rich aerobic media. The per-gene Table S1 is still served live at the Wisconsin mirror, which makes this the cleanest hash-pinnable per-gene essentiality table for MG1655. Note that the Database of Essential Genes lists 609 essential genes for this study and 296 for Baba 2006, both of which disagree with the source papers (620 and 303); use the paper tables. \\
\addlinespace[5pt]
46 & \textbf{Kochanowski 2017}\newline {\scriptsize Multi-omics campaign; promoter activities and absolute central metabolite concentrations on shared conditions} & \org{E. coli} & 4 & 2 wild type & 26 conditions & 26 & $3.7\times 10^{3}$ & reference-\allowbreak only & 95 fluorescent promoter reporters (an existing library expanded by 28, with 31 inactive promoters discarded) measured at steady state across 26 conditions spanning carbon sources, amino-acid supplementation and sub-lethal chloramphenicol, growth rates 0.1 to 1.5 per hour, in BW25113 plus a Crp deletion mutant. 47 central metabolites were quantified absolutely by ion-pairing UPLC-MS/MS in 23 of the 26 conditions, so dimensionality 142 is the union of 95 promoter activities and 47 metabolites, not a per-sample vector measured in every condition. No fluxes were measured. \\
\addlinespace[5pt]
47 & \textbf{Oyetunde 2019}\,\textsuperscript{\textbf{A}}\newline {\scriptsize Aggregation / support; titer (g/L), rate (g/L/h), yield (g/g)} & \org{E. coli} & 3 & 1,200 designs & 1 condition & $1.2\times 10^{3}$ & $3.6\times 10^{3}$ & reference-\allowbreak only & About 1200 experimentally realized E. coli cell factories hand-curated from about 100 papers covering more than 20 products, with features in six categories (carbon source, bioprocess conditions, genetic modifications, product properties, production metrics, undocumented factors) and TRY targets; data availability statement says all data are in the supplementary files. These are literature records, not measured strain-condition observations, so they must be resolved to primary source DOIs before counting as independent datasets. \\
\addlinespace[5pt]
48 & \textbf{Taniguchi 2010}\newline {\scriptsize Multi-omics campaign; single-cell protein copy number distribution and mRNA copy number per tagged gene} & \org{E. coli} & 3 & 1,018 designs & 1 condition & $1.0\times 10^{3}$ & $2.0\times 10^{3}$ & engineered-\allowbreak chassis & 1,018 confirmed strains of 1,400 attempted, each carrying a chromosomal C-terminal YFP fusion to one native gene at its native locus, imaged at about 4,000 cells per strain (96 strains per microfluidic device, 25 seconds each), with mRNA by RNA-seq and single-molecule FISH (137 genes had matched FISH). Per-strain dimensionality is only 2 summary values, so its value is the 1,018-gene tagged panel and the single-cell distributions, not depth per sample. The CGSC availability claim is unverified. \\
\addlinespace[5pt]
49 & \textbf{Haverkorn van Rijsewijk 2011}\newline {\scriptsize Metabolome / flux; 13C metabolic flux ratios in central carbon metabolism} & \org{E. coli} & 4 & 91 deletions & 2 conditions & 182$\dagger$ & $1.5\times 10^{3}$ & K-\allowbreak 12-\allowbreak KO & 91 regulator deletion mutants (81 transcription factors plus 10 sigma and anti-sigma factors) on glucose and galactose, giving 8 metabolic flux ratios per strain by GC-MS of proteinogenic amino acids with FiatFlux. Dimensionality is genuinely small; the value of this row is the 91-strain deletion axis, not depth. Do not confuse with the Nielsen commentary (DOI 10.1038/msb.2011.10, PMID 21451588). \\
\addlinespace[5pt]
50 & \textbf{Menasalvas 2025}\newline {\scriptsize Production campaign; isoprenol titer by GC-FID after biosensor-coupled pyrF growth selection} & \org{P. putida} & 3 & 165 deletions & 2 conditions & $1.3\times 10^{3}$$\dagger$ & $1.3\times 10^{3}$ & engineered-\allowbreak chassis & ATTRIBUTION CORRECTION: no author named Teng is on this paper. The Dryad DOI belongs to Menasalvas et al. with Thomas Eng as co-corresponding author, so the recalled Teng is almost certainly Eng. Over 165 deletion and overexpression strains across 70 gene loci, quadruplicate, titers at 24 and 48 h, reaching about 900 mg/L (36-fold). ONE study, four linked deposits, all confirmed: Dryad 10.5061/dryad.sbcc2frjq (2 files, 78.98 MB), PRIDE PXD061547 (An Isoprenol Biosensor for Combinatorial High Throughput Strain Engineering in Pseudomonas putida, 97 raw files with 8, 24 and 48 h timepoints), BioProject PRJNA1226229 (3 WGS runs), Zenodo 10.5281/zenodo.17155686 (AlphaFold output only). IMPORTANT on the CRISPRi library question: the paper does contain a pooled dCpf1 library of about 16,500 guides (3 per gene over 5,591 coding sequences), but there is NO released per-guide fitness or abundance matrix. Dryad holds only baseline per-guide read counts, a lost-guide list, per-round enriched-guide hit lists, and representative rather than per-replicate ONT reads, so a guide-by-sample matrix cannot be reconstructed. Proteomics covers about 2,500 to 3,000 proteins on a five-context subset, not every strain, so dimensionality is kept at 1. \\
\addlinespace[5pt]
\midrule \multicolumn{10}{@{}l}{\textbf{End of tranche 2. Rows 21--50 complete the fifty; rows below are the ranked reserve.}}\\ \midrule
51 & \textbf{Opgenorth 2019}\newline {\scriptsize Multi-omics campaign; 1-dodecanol titer (g/L) plus targeted proteomics of the engineered pathway proteins} & \org{E. coli} & 4 & 60 RBS variants & 2 conditions & 180$\ddagger$ & 900 & engineered-\allowbreak chassis+\allowbreak RBS & Two DBTL cycles over 60 engineered E. coli MG1655 strains; cycle 1 tested 36 strains varying ribosome-binding sites and reductase enzymes in one pathway operon, and cycle 2 designs from the machine-learning models raised titer 21\% to 0.83 g/L. The abstract states dodecanol plus all engineered-pathway protein concentrations were measured; the exact proteomics panel size and replicate count were not verified because the full text is paywalled, so dimensionality 5 and instances 180 are estimates. Conditions counted as 2 for the two DBTL cycles. \\
\addlinespace[5pt]
52 & \textbf{Valencia 2022 free fatty acids}\newline {\scriptsize Production campaign; free fatty acid and fatty acid methyl ester chain-length distribution at 48 h} & \org{P. putida} & 4 & 48 pathway designs & 3 conditions & 144$\ddagger$ & 864 & engineered-\allowbreak chassis & Included because its readout is genuinely multivariate: every sample yields a chain-length vector (C8, C10, C12:1, C14, C16, C16:1) rather than a scalar titer, so encoding it at dimensionality 1 would discard the study. Eight host backgrounds from deletions of three fatty acyl-CoA ligases, crossed with about six plasmid constructs carrying protein-engineered TesA thioesterase variants or MmFAMT, across three media including fourfold-diluted sorghum hydrolysate, n=3; maxima 670.9 mg/L total free fatty acid, 253.6 mg/L C8, 302.4 mg/L total methyl ester, 561.1 mg/L on hydrolysate. The genotype count is an estimate because the host by plasmid design is not fully crossed. Source data for Figures 2 to 6 sit in Supplementary Data 1, but data availability cites only the bare public-registry.jbei.org root with no folder number, so nothing is retrievable by accession. \\
\addlinespace[5pt]
53 & \textbf{Mohamed 2020 hydroxycinnamic TALE}\newline {\scriptsize Tolerance / robustness; growth rate and lag phase on p-coumarate and ferulate at increasing concentrations} & \org{P. putida} & 4 & 105 resequenced clones & 4 conditions & 420$\ddagger$ & 420 & evolved-\allowbreak WGS & Not on the recalled list but it is the best deposited P. putida tolerance evolution campaign. Six independent parallel replicates each for p-coumarate, ferulate and an equal-mass mixture, plus four replicates for a glucose control ALE, ramped to the solubility limit (40 g/L ferulate) or a 67-day cap; outcomes include a 37 h decrease in lag phase at 20 g/L p-coumarate and a 2.4-fold growth-rate increase at 30 g/L ferulate. PRJNA624660 confirmed by fetching the BioProject page: 102 SRA experiments and 105 BioSamples, with the ALEdb project named P. putida Hydroxycinnamic TALE. CAVEAT on genotype reconstruction: the deposit mixes endpoint populations with clonal isolates (1 to 2 intermediate isolates, the endpoint population, and a single endpoint isolate per condition), and only the isolates have a reconstructable single genotype; the population samples are blocked at the genotype level even though their reads are public. \\
\addlinespace[5pt]
54 & \textbf{Fang 2021 FFA CRISPRi}\newline {\scriptsize CRISPR library screen; free fatty acid titer (g/L) by GC} & \org{E. coli} & 4 & 108 guides & 1 condition & 324$\ddagger$ & 324 & K-\allowbreak 12+\allowbreak guide & Arrayed rather than pooled: 108 synthetic sgRNAs repressing 108 chromosomal genes across five metabolic modules, each strain measured for FFA titer by GC in three biological replicates, giving recoverable per-design phenotypes; 30 beneficial genes from the screen plus 26 more from omics, 56 total, and 30.0 g/L in fed-batch. Accessions are named in the paper but were not fetched. \\
\addlinespace[5pt]
55 & \textbf{Jones 2015 ePathOptimize}\newline {\scriptsize Combinatorial design; violacein titer (mg/L) by HPLC} & \org{E. coli} & 4 & 307 promoter variants & 1 condition & 307 & 307 & engineered-\allowbreak chassis+\allowbreak promoter & Five IPTG-inducible mutant T7 promoters of graded strength were placed on each of the five violacein pathway genes vioABCDE in E. coli BL21star(DE3), a theoretical 3125-member library; 107 colonies were screened first (3.4\% of the library) and a second enriched library of 200 mutants was screened, giving 307 genotype-linked titers. Best screen hit 238 mg/L, 63-fold over control, rising to 1829 plus or minus 46 mg/L after fermentation optimization. Deoxyviolacein and pathway intermediates were also quantified for selected strains, so the per-instance panel may exceed 1. \\
\addlinespace[5pt]
56 & \textbf{Jervis 2019}\newline {\scriptsize Combinatorial design; limonene titer (mg per L organic phase)} & \org{E. coli} & 4 & 88 RBS variants & 1 condition & 264$\ddagger$ & 264 & engineered-\allowbreak chassis+\allowbreak RBS & Two RBS libraries in E. coli DH10-beta: pGLlib pairs 12 designed RBS variants each for trAg-gpps and trMs-limS, 144 possible combinations, from which 360 colonies were screened in deep-well plates and 64 clones sequenced, covering 56 of the 144 combinations; a second 5184-member pMVA2 library had 156 colonies screened and a designed 32-combination reduced library built and tested. Genotypes 88 counts only sequence-resolved designs (56 plus 32), since unsequenced colonies have titers without a recoverable genotype; instances assumes triplicate. \\
\addlinespace[5pt]
57 & \textbf{Kang 2026 isoprenyl acetate}\newline {\scriptsize Production campaign; isoprenyl acetate and isoprenol titer from mixed glucose and xylose} & \org{P. putida} & 4 & 25 pathway designs & 3 conditions & 225$\ddagger$ & 225 & engineered-\allowbreak chassis & Found in a ProteomeXchange sweep rather than the recalled list, and it belongs only if the downstream ester is admitted: the product is isoprenyl acetate, an ATF1-esterified derivative of isoprenol, at 1.9 g/L in fed-batch with a yield of 0.067 g per g total sugar. PXD067010 confirmed by fetching its PRIDE record, which names this DOI and PMID and lists Petzold as lab head. Engineering combines a heterologous alcohol acetyltransferase, esterase deletions to stop product degradation, glucose-xylose co-utilization, and reinforced acetyl-CoA flux, across tube, flask and 2 L fed-batch in biological triplicate. Strain and instance counts are estimates; the deposit does not state a run count. Shares Carruthers and Lee with the flagship isoprenol rows, so deduplicate genotypes. \\
\addlinespace[5pt]
58 & \textbf{Schmidt 2025 3-hydroxyacid PKS}\newline {\scriptsize Production campaign; LC-MS titer of three 3-hydroxyacids in supernatant at 48 h} & \org{P. putida} & 4 & 12 pathway designs & 2 conditions & 72$\ddagger$ & 216 & engineered-\allowbreak chassis & This is the modular type I PKS row. JBEI public registry folder 887 (Schmidt et al 2025) was confirmed through the REST endpoint: 58 STRAIN entries with public read access, which is an upper bound on strains built and the reason the genotype count is a figure-level estimate rather than a stated total. Three products quantified (3H24DMPA, 3H4MPA at 18.2 mg/L, 3H4MHA at 6.4 mg/L), n=3 for titers, up to seven media variants that are not fully crossed, single 48 h harvest. The genotype classes here exceed the usual four and need extra encoding: alongside deletions and plasmid PKS expression there are mini-Tn7 and BxB1 chromosomal integrations, C-terminal ssrA degron tags on FabD (LAA active, DAS intermediate, LDD inert), and acyltransferase-domain swaps in the LipPKS1 loading module. Accompanying RB-TnSeq is said to be at fit.genomics.lbl.gov, unconfirmed. A sibling deposit from the same group, PXD069956 (Engineering an Extremely Hybrid PKS for Adipic Acid Production, Klass et al., ACS Synth Biol, DOI 10.1021/acssynbio.5c00972, PMID 42339614), covers the same PKS platform in P. putida and E. coli if the modality is worth widening. \\
\addlinespace[5pt]
59 & \textbf{Alper 2005 lycopene}\newline {\scriptsize Combinatorial design; lycopene content per cell dry weight} & \org{E. coli} & 4 & 64 designs & 1 condition & 192$\ddagger$ & 192 & K-\allowbreak 12-\allowbreak KO & Two independent target-discovery routes, a genome-scale-model systematic search and a transposon-insertion combinatorial search, produced two disjoint gene sets whose exhaustive combination gave a defined set of 64 knockout strains, the best 8.5-fold over recombinant K-12 wild type and 2-fold over the engineered parent. The deletion set is in a K-12 background that also carries the heterologous lycopene pathway; transposon insertions were used for discovery, not as the final genotypes. Replicate counts were not verified. \\
\addlinespace[5pt]
60 & \textbf{ActiveOpt valine (Kumar 2021)}\newline {\scriptsize Combinatorial design; valine elemental carbon yield (\% of maximum theoretical yield from glucose plus acetate)} & \org{E. coli} & 4 & 93 RBS variants & 1 condition & 186$\ddagger$ & 186 & engineered-\allowbreak chassis+\allowbreak RBS & New experimental dataset reported in this paper: 39 plasmids expressing ilvBN*DE and ilvIH*C/C*-ygaZH with varied RBS strengths, tested as 89 pairwise plasmid combinations in strain PYR003a (BW25113 delta-aceE delta-gdhA delta-poxB delta-ldhA delta-recA), plus 4 new combinations chosen by ActiveOpt, so 93 genotype-linked yields; best strain 45\% elemental carbon yield, 54.7\% of maximum theoretical. A minimum of two biological replicate colonies per experiment, 48 h shake flask endpoint, no time course. \\
\addlinespace[5pt]
61 & \textbf{Banerjee 2024}\newline {\scriptsize Production campaign; isoprenol titer plus growth and glucose consumption, shake flask and fed-batch} & \org{P. putida} & 4 & 19 deletions & 3 conditions & 171$\ddagger$ & 171 & KT2440-\allowbreak KO & The recalled 3.5 g/L fed-batch figure is confirmed in the abstract, against 1.1 g/L in batch and over tenfold above the starting strain. Eight deletion targets came from bilevel optimization and constrained minimal cut sets on iJN1462, and the best strain carries five deletions. The strain count is an estimate read from the preprint because the published version is paywalled. Proteomics here is targeted MRM of pathway enzymes only, not a discovery panel, so dimensionality stays 1. Linked assets: PXD039868 on PanoramaPublic plus six Experimental Data Depot studies at edd.jbei.org, which were not opened. \\
\addlinespace[5pt]
62 & \textbf{Wang 2018 phosphatase CRISPRi}\newline {\scriptsize CRISPR library screen; lycopene content (and beta-carotene in follow-up strains)} & \org{E. coli} & 4 & 56 guides & 1 condition & 168$\ddagger$ & 168 & engineered-\allowbreak chassis & A CRISPRi knockdown library over 56 phosphatase-encoding genes built in a lycopene overproducer; 28 of the 56 knockdowns impaired lycopene synthesis, and combinatorial knockdowns plus knockouts were then tested in lycopene and beta-carotene overproducers. The gene set is small and fully named, so per-design values should be recoverable from the figures, but replicate counts were not verified so instances is an estimate. \\
\addlinespace[5pt]
63 & \textbf{Jain 2015 1,2-propanediol}\newline {\scriptsize Production campaign; 1,2-propanediol titer (g/L) and yield (g/g glucose)} & \org{E. coli} & 4 & 20 designs & 1 condition & 60$\ddagger$ & 120 & engineered-\allowbreak chassis & I could not confirm the existence of a distinct growth-coupled 1,2-propanediol design-validation campaign: repeated searches surfaced only computational growth-coupling frameworks (OptKnock, minimal cut sets, gcFront) plus separate 1,2-PDO engineering papers, so the 1,2-propanediol row is filled by the best-documented reconstructable E. coli campaign, which reached 5.13 g/L and 0.48 g/g glucose via an optimal minimal enzyme set, flux channeling, raised NADH availability and improved anaerobic growth. The citation is sourced; the genotype and instance counts are unverified estimates, hence confidence recall. \\
\addlinespace[5pt]
64 & \textbf{Lim 2025 isoprenol TALE}\newline {\scriptsize Tolerance / robustness; maximum specific growth rate under isoprenol challenge plus isoprenol titer with the pIY670 plasmid} & \org{P. putida} & 4 & 46 resequenced clones & 2 conditions & 114$\dagger$ & 114 & evolved-\allowbreak WGS & The recalled design is correct and reported: WT, IPL300 and IPL400 each evolved in four independent lineages (12 isoprenol lineages), plus four extra WT lineages tolerized against formaldehyde as controls, so 16 lineages total, isoprenol ramped 4 to 8.5 g/L over 143 to 353 generations. Only the 46 whole-genome resequenced strains (16 endpoint clones plus about 30 intermediates) have a reconstructable genotype; the evolved populations themselves do not and are blocked at the genotype level. 158 unique mutations across 73 regions; proteomics quantified an average of 2,375 of 5,565 putative proteins. Linked assets: PXD054609 (confirmed, citing this PMID), GEO GSE281392, and ALEdb project Pputida\_isoprenol\_TALE. IMPORTANT: the stated resequencing accession PRJNA1187681 does NOT resolve at NCBI BioProject, NCBI SRA or ENA as of 2026-09-24, so the raw reads are not retrievable; the mutation calls survive only in Supplementary Dataset 1. \\
\addlinespace[5pt]
65 & \textbf{Farasat 2014 neurosporene (ActiveOpt case 2)}\newline {\scriptsize Combinatorial design; specific neurosporene productivity (microgram per gCDW per hour)} & \org{E. coli} & 4 & 101 RBS variants & 1 condition & 101 & 101 & engineered-\allowbreak chassis+\allowbreak RBS & Kept as a row separate from the valine campaign because the design space is independent (crtEBI RBS library, different host and readout). The neurosporene measurements are Farasat et al. 2014's, not new work: 73 designed crtEBI expression constructs measured as exploration experiments plus 28 kinetic-model-designed extrapolation constructs, productivity rising from 196.3 to 286 microgram per gCDW per hour. ActiveOpt (Kumar 2021, doi 10.1016/j.ymben.2021.06.009) is a retrospective reanalysis of that dataset and generated no new neurosporene strains, so the primary citation is Farasat 2014. \\
\addlinespace[5pt]
66 & \textbf{Lim 2020 ionic liquid ALE}\,\textsuperscript{\textbf{B}}\newline {\scriptsize Tolerance / robustness; growth performance at high protic ionic liquid concentrations} & \org{P. putida} & 4 & 10 resequenced clones & 3 conditions & 90$\ddagger$ & 90 & evolved-\allowbreak WGS & Bibliography confirmed through CrossRef (Green Chemistry 22(17):5677-5690); the journal is not indexed in PubMed, so the PMID is deliberately left empty rather than guessed. Evolved strains grow at up to 4 percent triethanolammonium acetate and 8 percent triethylammonium hydrogen sulfate where the wild type cannot, with mutations in relA, gacS, oprB, fleQ, tktA and uvrY for minimal media and PP\_5350, emrE, oprD and PP\_5324 for the ionic-liquid-specific arms; PP\_5350 (an RpiR-family regulator upregulating the glyoxylate cycle) and the emrE efflux pump were validated by reverse engineering. Marked BLOCKED because no sequencing deposit could be located: the publisher page returns HTTP 403, OSTI dropped the connection, and no BioProject or SRA accession surfaced anywhere, so the evolved genotypes are not retrievable even though the mutations are named in the text. Strain and instance counts are estimates. \\
\addlinespace[5pt]
67 & \textbf{Ling 2022 muconate}\newline {\scriptsize Production campaign; cis,cis-muconate titer, molar yield and volumetric rate on glucose and xylose} & \org{P. putida} & 4 & 11 deletions & 4 conditions & 66$\ddagger$ & 66 & engineered-\allowbreak chassis & The best documented muconate campaign with per-strain titers: 11 strains carry titer or growth measurements and Table 1 lists 24 strains, across four carbon conditions in flask, plate reader and 0.5 L fed-batch, reaching 33.7 g/L at 0.18 g/L/h and 46 percent molar yield. PRJNA783062 confirmed by fetching the BioProject page, titled Sequencing of 5 ALE isolates on xylose (NREL, 2021-11-23), so only the genome resequencing is deposited and there is no proteomics accession. Genotype encoding is unusually mixed and does not fit a single modality: rational deletions and plasmid pathways sit alongside ALE-derived promoter mutations in the PP\_2569 promoter, xylE point mutations (A62V, A455V) and a roughly 227.8 kb genome duplication from PP\_5050 to PP\_5242, none of which are designed edits. \\
\addlinespace[5pt]
68 & \textbf{Werner 2023 beta-ketoadipate}\newline {\scriptsize Production campaign; beta-ketoadipate titer, productivity and molar yield from lignin-derived aromatics} & \org{P. putida} & 4 & 8 deletions & 4 conditions & 64$\ddagger$ & 64 & KT2440-\allowbreak KO & Eight strains with titers across four aromatic feeds (p-coumarate, ferulate, a 3:1 molar mixture, and corn stover alkaline pretreated liquor extractives), reaching 44.5 g/L at 1.15 g/L/h from model compounds and 25 g/L at 0.66 g/L/h from real corn stover aromatics; genotypes are deletions of pcaIJ, crc, pobAR, lvaE, gacA and gacS plus chromosomal overexpression at the fpvA locus. It has the best per-strain numerical table of this cluster, with every main-text data point tabulated in data S1. But there is no repository accession at all and strains are available only from NREL under a material transfer agreement, so accession\_confirmed is false because no accession exists. One arm is informative as a negative: the gacA variants underperformed. \\
\addlinespace[5pt]
69 & \textbf{de Siqueira 2025}\newline {\scriptsize Tolerance / robustness; growth rate on glucose-acetate mixtures plus isoprenol titer by GC-FID} & \org{P. putida} & 4 & 7 resequenced clones & 3 conditions & 63$\dagger$ & 63 & evolved-\allowbreak WGS & Both recalled accessions verified by fetching the repository records. PXD055153 is titled Restoration of growth and isoprenol production in dual carbon source media by an acetate-sensitive Pseudomonas putida strain, submitted 2024-08-23, linked to this DOI. PRJNA1153078 (Pseudomonas putida KT2440 tolerization on acetate and isoprenol) holds exactly 5 SRA experiments and 5 BioSamples, matching the 5 resequenced isolates, four tolerized Sigma-class strains plus the parental producer; 2 of the 7 characterized strains were not resequenced and are blocked at the genotype level. ONE study with linked assets PXD055153 plus PRJNA1153078. Proteome panel size is never stated numerically, so dimensionality is kept at 1 rather than guessed. \\
\addlinespace[5pt]
70 & \textbf{Yunus 2026}\newline {\scriptsize Production campaign; isoprenol titer under predicted CRISPRi knockdowns} & \org{P. putida} & 4 & 20 guides & 1 condition & 60$\ddagger$ & 60 & KT2440+\allowbreak guide & Found while sweeping for CRISPRi isoprenol work and it clearly belongs: FluxRETAP and VAMMPIRE were used to pick downregulation targets and to assemble arrays of up to five sgRNAs, reaching nearly 1.5 g/L isoprenol by knocking down PP\_4118 (alpha-ketoglutarate dehydrogenase). Bibliography verified through NCBI eutils. Strain and instance counts are estimates: the abstract does not state them and the full text is paywalled. No accession is given in the indexed metadata, so data\_location is deliberately empty rather than guessed. Shares authors and the pIY670-era pathway with Carruthers 2025 and Banerjee 2024, so genotype deduplication against those rows is required. \\
\addlinespace[5pt]
71 & \textbf{Banerjee 2020 indigoidine}\newline {\scriptsize Production campaign; indigoidine titer, rate and yield plus growth, scaled from deep-well plate to 2 L bioreactor} & \org{P. putida} & 4 & 2 guides & 8 conditions & 48$\ddagger$ & 48 & KT2440+\allowbreak guide & This is the growth-coupled indigoidine CRISPRi campaign, and it needs a CORRECTION to the recalled framing. The counts are right but the object is not: 63 constrained-minimal-cut-set solutions were analyzed, one feasible cut set hit 14 metabolic reactions, which mapped through GPRs to 16 single-copy genes, of which 2 (mqo-I and cynT) were dropped as essential by genome-wide RB-TnSeq, leaving 14 genes knocked down. Those 14 guides sit in ONE multiplexed dCpf1 gRNA array in a SINGLE strain. There are not 14 CRISPRi strains, so the genotype axis here is 2 (engineered versus empty-vector control), not 14, and anything that ingests this as a 14-strain panel is wrong. Conditions are 2 carbon sources crossed with 4 culture formats. RNA-seq is genome-wide (about 5,500 genes, estimated from the annotation, not stated); the proteomics is targeted SRM, so it is a selected reaction list rather than a discovery panel. BioProjects PRJNA580539 through PRJNA580574 (both endpoints resolve; the paper typos the upper bound); proteomics at PanoramaWeb, not opened. \\
\addlinespace[5pt]
72 & \textbf{Wang 2023 BATCH biosensor CRISPRi}\,\textsuperscript{\textbf{B}}\newline {\scriptsize CRISPR library screen; biosensor fluorescence sort, then titer (mg/L) on selected sgRNA variants} & \org{E. coli} & 4 & 20 guides & 2 conditions & 40$\ddagger$ & 40 & engineered-\allowbreak chassis & Doubly mismatched sgRNA pools give graded knockdown of 20 central-carbon-metabolism genes, screened with a PadR-based p-coumaric acid biosensor and an HpdR-based butyrate biosensor: pfkA plus ptsI variants raised p-coumarate 40.6\% to 1308.6 mg/L from glycerol, and sucA or ldhA variants raised butyrate 19.0\% and 25.2\%. The total number of mismatched sgRNA variants in the pools was not confirmed, and the pooled sort yields no per-variant phenotype except for the few reconstructed winners, hence blocked. Conditions 2 counts the two product-biosensor campaigns. \\
\addlinespace[5pt]
73 & \textbf{Czajka 2022}\newline {\scriptsize Production campaign; indigoidine titer by A612 in DMSO plus growth rate} & \org{P. putida} & 4 & 6 guides & 1 condition & 36$\dagger$ & 36 & KT2440+\allowbreak guide & The direct follow-up to Banerjee 2020 on the same 14-target multiplexed CRISPRi strain, adding an optimized Cpf1 ribosome binding site for a 1.6-fold gain, PHA-operon deletion strains, RNA-seq, targeted proteomics and 13C metabolic flux analysis. Six strains carry titers in one condition (M9, 10 g/L glucose, arabinose plus IPTG), n at least 6 for production profiling, sampled at 24, 48, 72 and 120 h. Kept as its own row rather than folded into Banerjee 2020 because the genotypes and the readout differ. Data availability names only public-registry.jbei.org with no folder number and states explicitly that everything else is in the manuscript and supplementary files, so there is no GEO, SRA, PRIDE or MassIVE deposit and nothing is machine-retrievable. \\
\addlinespace[5pt]
74 & \textbf{Fong 2005 growth-coupled lactate}\newline {\scriptsize Production campaign; lactate titer (g/L), growth rate, lactate secretion rate} & \org{E. coli} & 4 & 3 deletions & 1 condition & 11 & 33 & K-\allowbreak 12-\allowbreak KO & Three OptKnock-designed knockout strains were built and adaptively evolved to yield 11 evolved lactate-producing strains, growing on 2 g/L glucose at 37 C with titers 0.87 to 1.75 g/L and secretion rates coupled to growth rate; computational post-evolution growth-rate predictions were within 10\% in 38 of 50 cases. Only 3 distinct engineered genotypes exist, so this is a small validation campaign rather than a library. \\
\addlinespace[5pt]
75 & \textbf{Yu 2016 PHBA}\newline {\scriptsize Production campaign; para-hydroxybenzoic acid titer from glucose via chorismate, plus carbon yield} & \org{P. putida} & 4 & 7 deletions & 1 condition & 21$\dagger$ & 21 & KT2440-\allowbreak KO & The recalled six strains plus wild type is exact: S0 wild-type control and S1 through S6, built from markerless deletions of pobA, pheA, trpE and hexR combined with plasmid-borne E. coli ubiC and feedback-resistant aroG-D146N; S6 reached 1.73 g/L at 18.1 percent C-mol per C-mol. One defined medium, one 44 h endpoint, biological triplicates, which makes it the cheapest row here to ingest. A single S6 fed-batch run adds a short time series. Up to eight co-measured scalars exist per instance (PHBA, biomass, glucose, acetate, alpha-ketogluconate, pyruvate, lactate, succinate) if each counts as a phenotype. No proteomics, no transcriptomics, no repository deposit of any kind, so accession\_confirmed is false because no accession exists to confirm. \\
\addlinespace[5pt]
\end{longtable}
\endgroup
\end{landscape}
